\section*{Цель работы}
\begin{quote}
    Исследование метрологических характеристик цифровых приборов, их применение для измерения физических величин и оценка погрешностей результатов измерений.
\end{quote}

\section*{Задачи}
\begin{enumerate}
    \item Ознакомиться с инструкцией по применению исследуемого цифрового измерительного прибора (ЦИП).
    \item Определить шаг квантования (квант) исследуемого ЦИП в режиме омметра для различных (по указанию преподавателя) пределов измерения.
    \item Экспериментально определить метрологические характеристики ЦИП в режиме омметра:
    \begin{itemize}
        \item статическую характеристику преобразования и построить ее график $R_\text{п} = F(R)$;
        \item инструментальную погрешность по всему диапазону измерений для выбранного предела;
        \item построить график инструментальной погрешности и определить ее аддитивную и мультипликативную составляющие.
    \end{itemize}
    \item Измерить сопротивления ряда резисторов и оценить основную погрешность результатов измерения.
\end{enumerate}

\section*{Теоретические сведения}
В цифровых измерительных приборах (ЦИП) результаты измерений представляются в цифровом виде. В отличие от аналоговых приборов, показания ЦИП изменяются дискретно на единицу младшего разряда (ЕМР).

\textit{Статическая характеристика преобразования} устанавливает связь между преобразуемой входной величиной $x$ и результатом преобразования $x_\text{п}$ (показаниями ЦИП). Результат может принимать только квантованные значения:
\begin{equation*}
    x_\text{п} = Nq
\end{equation*}
где $N$ — десятичное целое число, а $q$ — шаг квантования (квант) величины $x$. Статическая характеристика имеет ступенчатый вид, как показано на \cref{fig:static_char}. Размер ступени равен кванту $q$.

\begin{figure}[H]
    \centering
    \includegraphics[width=0.4\linewidth]{figs/static_char.pdf}
    \caption{Идеальная (пунктир) и реальная (сплошная линия) статические характеристики ЦИП}
    \label{fig:static_char}
\end{figure}

Значение кванта $q$ связано с пределом измерений $x_\text{max}$ и максимальным числом уровней квантования $N_\text{max}$ соотношением:
\begin{equation*}
    q = \frac{x_\text{max}}{N_\text{max}}
\end{equation*}

Методическая погрешность ЦИП, или \textit{погрешность квантования}, обусловлена дискретностью представления информации и принимается равной половине кванта, т.е. $\pm 0.5q$.

В реальном ЦИП статическая характеристика смещается из-за наличия инструментальных погрешностей. Смена показаний происходит при значениях входной величины $x$, отличных от идеальных значений $Nq$. Абсолютная инструментальная погрешность для конкретного показания $x_\text{п}$ определяется по отличию реальной характеристики от идеальной:
\begin{equation}
    \Delta x_\text{и} = x_\text{п} - 0.5q - x
    \label{eq:instr_error}
\end{equation}
где $x$ — значение входной величины, при котором происходит смена показаний на $x_\text{п}$.

Погрешность ЦИП часто моделируется линейной зависимостью:
\begin{equation*}
    \Delta x = a + bx
\end{equation*}
где $a$ — \textit{аддитивная} составляющая погрешности (не зависит от значения измеряемой величины), а $bx$ — \textit{мультипликативная} составляющая (пропорциональна измеряемой величине).

\section{Порядок выполнения работы}

\subsection{Подготовка к работе}
\begin{enumerate}
    \item Подключите магазин сопротивлений ко входам ЦИП, работающего в режиме омметра.
    \item Выберите предел измерения согласно указанию преподавателя. Определите для этого предела значение единицы младшего разряда (кванта) $q$.
    \item Убедитесь, что шаг изменения сопротивления магазина $q_\text{м}$ значительно меньше кванта прибора $q$ ($q \gg q_\text{м}$).
\end{enumerate}

\subsection{Определение начального участка статической характеристики}
\begin{enumerate}
    \item Установите на магазине сопротивлений нулевое значение $R=0$.
    \item Плавно увеличивайте сопротивление магазина $R$, следя за показаниями ЦИП $R_\text{п}$.
    \item В момент, когда показания $R_\text{п}$ изменятся на единицу младшего разряда (например, с $0.001~\text{кОм}$ на $0.002~\text{кОм}$), зафиксируйте новое показание $R_\text{п}$ и соответствующее ему значение сопротивления магазина $R$.
    \item Занесите пару значений ($R_\text{п}$, $R$) в \cref{tab:static_char_protocol}.
    \item Повторите процедуру для 8--9 последовательных значений $R_\text{п}$.
\end{enumerate}

\subsection{Определение инструментальной погрешности}
\begin{enumerate}
    \item Выберите 8--10 точек, равномерно распределенных по всему выбранному диапазону измерений.
    \item Для каждой точки установите на магазине сопротивлений соответствующее значение $R$.
    \item Снимите показание ЦИП $R_\text{п}$.
    \item Занесите пары значений ($R_\text{п}$, $R$) в \cref{tab:instr_err_protocol}.
\end{enumerate}

\subsection{Измерение сопротивлений}
\begin{enumerate}
    \item Получите у преподавателя набор резисторов для измерения.
    \item Для каждого резистора выберите на ЦИП наиболее подходящий диапазон измерения (обеспечивающий максимальную точность).
    \item Измерьте сопротивление каждого резистора.
    \item Занесите в \cref{tab:resistors_protocol} номер резистора, выбранный диапазон и показание прибора $R_\text{п}$.
\end{enumerate}

\section{Обработка результатов}
\begin{enumerate}
    \item По данным \cref{tab:static_char_protocol} построить график начального участка статической характеристики $R_\text{п} = F(R)$.
    \item Для каждой точки из \cref{tab:instr_err_protocol} рассчитать абсолютную инструментальную погрешность $\Delta R_\text{и}$ по формуле \cref{eq:instr_error}.
    \item Построить график зависимости инструментальной погрешности от измеряемого сопротивления $\Delta R_\text{и} = F(R)$.
    \item По графику $\Delta R_\text{и} = F(R)$ определить аддитивную ($a$) и мультипликативную ($b$) составляющие погрешности.
    \item Для каждого резистора из \cref{tab:resistors_protocol} рассчитать абсолютную и относительную погрешности измерения. Представить результат в виде $R = R_\text{п} \pm \Delta R$.
    \item Сформулировать выводы о характере изменения погрешности ЦИП и дать рекомендации по выбору предела измерения.
\end{enumerate}

%% \cref{fig:static_char} соответствует Рис. 3.1 из методического пособия.